\subsection{Pivots, Back Substitution, and Row Echelon Form}

A pivot is a leading nonzero entry used to eliminate other entries in its column. Pivots mark independent information and identify the variables that can be solved.

\begin{examplebox}[Back substitution in a $2\times2$ system]
For
\[
\left[\begin{array}{cc|c}
1&2&5\\
0&3&6
\end{array}\right],
\]
the second row gives $y=2$, and the first row then gives $x=1$.
\end{examplebox}

\begin{definitionbox}[Row echelon form]
A matrix is in row echelon form when leading nonzero entries move strictly to the right as one moves downward, and all zero rows are placed at the bottom.
\end{definitionbox}

A row such as
\[
\left[\begin{array}{cc|c}0&0&0\end{array}\right]
\]
represents $0=0$ and imposes no new restriction.